\documentclass[a4paper,12pt,twoside, titlepage]{book}
\usepackage[latin,french]{babel}

\usepackage{fontspec}
\defaultfontfeatures[Old Standard]{%utiliser oldstandard pour imprimer le grec
Extension = .otf, 
UprightFont = OldStandard-Regular,  
ItalicFont = OldStandard-Italic,  
BoldFont = OldStandard-Bold,  
BoldItalicFont = OldStandard-BoldItalic,  
SlantedFont = OldStandard-Regular,
SlantedFeatures = {FakeSlant=0.25},
BoldSlantedFont = OldStandard-Bold,
BoldSlantedFeatures = {FakeSlant=0.25}
}
\setmainfont{Old Standard}
\usepackage[poetry=verse]{ekdosis}

\SetAlignment{
texts=translation;edition,
apparatus=edition
}

\title{exercice Ekdosis}
\author{Clara Grometto}
\date{19/04/2026}

\AtBeginEnvironment{edition}{\selectlanguage{latin}} %important pour la ponctuation et les hyphénations. En théorie, combiné à Babel, on pourrait enlever les ~ pour les espaces insécables
\AtBeginEnvironment{translation}{\selectlanguage{french}}


\DeclareWitness{M}{\emph{M}}{Codex Mediceus}
[origDate={saec. \textsc{v}}]
\DeclareWitness{P}{\emph{P}}{Codex Palatinus}
  [origDate={saec. \textsc{iv}}]
\DeclareWitness{R}{\emph{R}}{Codex Romanus}
  [origDate={saec. \textsc{v}}]
\DeclareWitness{V}{\emph{V}}{Schedae Veronenses rescriptae}
  [origDate={saec. \textsc{v}}]

\DeclareShorthand{coddPR}{\emph{codd.}}{P,R}
  
\DeclareWitness{a}{\emph{α}}{Codex Mediolanensis biblioth. Ambrosianae, olim Petrarcae} %les id en alphabet latin car c'est trop compliqué de les passer en grec
  [origDate={saec. \textsc{vxiii-xiv}}]
\DeclareWitness{c}{\emph{γ}}{Codex Guelferbytanus Gudianus 2\textsuperscript{o}, 70.}
  [origDate={saec. \textsc{ix}}]

  \DeclareSource{Serv}{\textsc{Serv.}} %@Servius (4e s. apr.) a commenté les Bucoliques de Virgile (https://www.perseus.tufts.edu/hopper/text?doc=Perseus%3atext%3a2007.01.0092). Donc pour établir le texte des Buc. de Virg., on consulte aussi ce genre de textes anciens. C'est une source. 
  \DeclareSource{Daniel}{\textsc{Daniel.}} 
  \DeclareSource{Philarg}{\textsc{Philarg.}} %https://fr.wikipedia.org/wiki/Junius_Philargyrius
  \DeclareSource{Donat}{\textsc{Donat.}} 
  \DeclareSource{Schol}{\textsc{Schol. Bern.}}

  \DeclareSource{Quintil}{\textsc{Quintil.}} %Quintil et Asper ne sont pas recensés dans le CS mais ils pourraient être considérés comme des sources. 
  \DeclareSource{Asper}{\textsc{Asper.}}

\SetLineation{lineation=none,vmodulo} %numérotation toutes les 5 lignes, plus difficile à implémenter pour la version française. 

\DeclareApparatus{default}

\begin{document}


\maketitle



\begin{alignment}

\begin{translation}

\begin{center}\huge\textsc{Virgile}\end{center}

\begin{center}\large\textsc{Bucoliques}\end{center}

\begin{center}\large I\end{center}

\end{translation}
\begin{edition}

\begin{center}\huge\textsc{P. Vergili Maronis}\end{center}

\begin{center}\large\textsc{Bvcolica}\end{center}

\begin{center}\large I\end{center}
\end{edition}

\end{alignment}


\begin{alignment}

%%% --- PAR. 1 --- %%%

\begin{translation}

\begin{center}\textsc{Mélibée}\end{center}

Toi, Tityre, étendu sous le couvert d'un large hêtre, tu essaies un air silvestre sur un mince pipeau; nous autres, nou quittons notre pays et nos chères campagnes; loin du pays nous sommes exilés; toi, Tityre, nonchalant sous l'ombrage, tu apprends aux bois à redire le nom de la belle Amaryllis. 

\end{translation}

\begin{edition}

\begin{center}\textsc{Meliboevs}\end{center}
\begin{ekdverse}
Tityre, tu patulae recubans sub tegmine fagi\\
\app{
  \lem[wit={coddPR}]{siluestrem}
  \rdg[source={Quintil}]{agrestem}
  \note{(IX, 4, 85) cf. B. 6, 8}
  }
tenui musa meditaris auena~;\\
nos patriae finis et dulcia linquimus arua~;\\
nos patriam fugimus~; tu, Tityre, lentus in umbra,\\
\app{
  \lem[wit={P,R,a,c}]{formosam}
  \rdg[source={Asper}, alt={---onsam}]{formonsam}
  \note{(\emph{Append. Seruiana, Hagen}, p.~540)}
} resonare doces Amaryllida siluas.\\
\end{ekdverse}

\end{edition}

%%% --- PAR. 2 --- %%%

\begin{translation}

\begin{center}\textsc{Tityre}\end{center}


O Mélibée, c'est à un dieu que nous devons ces loisirs~; car Il sera pour moi toujours un dieu~; son autel, une tendre victime, un agneau de nos bergeries, souvent l'ensanglantera. Grâce à Lui, mes génisses ont le droit de paître en liberté, comme tu vois, et moi-même celui de jouer mes airs préférés sur un roseau rustique.

\end{translation}

\begin{edition}
    \begin{center}\textsc{Tityrvs}\end{center}
    \begin{ekdverse}
O Meliboee, deus nobis haec otia fecit~:\\
namque erit ille mihi semper deus~; illius aram\\
saepe tener nostris ab ouilibus imbuet agnus.\\
Ille meas errare boues, ut cernis, et ipsum\\
ludere quae uellem calamo permisit agresti.\\
\end{ekdverse}
\end{edition}


%%% --- PAR. 3 --- %%%


\begin{translation}

\begin{center}\textsc{Mélibée}\end{center}


Je ne suis point jaloux, mais étonné plutôt~; partout, dans les campagnes, il y a tant de désordre~! Vois~: mes pauvres chèvres, je les pousse, dolent, droit devant moi~; et celle-ci, Tityre, j'ai peine à la tirer~: ici, dans une touffe de coudriers, elle vient de mettre bas deux bessons, espoir du troupeau, hélas~! laissés sur le roc nu.

\end{translation}
\begin{edition}

\begin{center}\textsc{Meliboevs}\end{center}
\begin{ekdverse}
Non equidem inuideo, miror magis~: undique totis\\
usque adeo \app{
  \lem[wit={a}, source={Quintil}]{turbatur}
  \note{(I, 4, 28) \emph{probat} \textsc{Serv.}} %(Servius approuve la leçon turbatur donnée par Quintilien)
  \rdg[wit={P,R,c}, alt={---mur}]{turbamur} %Attention, ici ce qui va entre accolade c'est ce qui est envoyé éventuellement en XML ; en XML, on veut ce qui se trouve dans le mss, soit turbamur et non pas ce qui se trouve sous forme abrégée dans l'apparat 
} agris~! En ipse capellas\\
\app{
  \lem[wit={P,R}]{protinus}
  \rdg[wit={a,c}, alt={---tenus}]{protenus}
} aeger ago~; hanc etiam uix, Tityre, duco~:\\
hic inter densas corylos modo namque gemellos,\\
spem gregis, a~! silice in nuda conixa reliquit.\\
\end{ekdverse}

\end{edition}
\end{alignment}

\end{document}



